\chapter{Fourier 級数}
    \section{定義と背景の整理}
        $T>0$ 周期関数 $f:\realNumbers\to\realNumbers$ に対して，次式を $f$ の Fourier 級数と呼ぶ。
        \[ a_0/2 + \sum_{d=1}^\infty \bracks*{a_d\cos\parens*{d\frac{t}{T}2\pi} + b_d\cos\parens*{d\frac{t}{T}2\pi}} \]
        ここに $a_0,a_d,b_d\;(d=1,2,3,\dots)$ は次式で与えられる。
        \[ a_0 = \frac{1}{T/2}\integrate{-T/2}{T/2}{f(t)}{}{t},\quad a_d = \frac{1}{T/2}\integrate{-T/2}{T/2}{f(t)\cos\parens*{d\frac{t}{T}2\pi}}{}{t},\quad b_d = \frac{1}{T/2}\integrate{-T/2}{T/2}{f(t)\sin\parens*{d\frac{t}{T}2\pi}}{}{t} \]
        \subsection{無限次元ベクトル空間との関係}
            区間 $[-T/2,T/2]$ 上の Riemann 可積分な実数値関数全体の成すベクトル空間 $V$ 上の内積 $\inprod{\cdot}{\cdot}:V\times V\to\realNumbers$ を次式で定義する。
            \[ \inprod{f}{g} \coloneq \frac{1}{T/2}\integrate{-T/2}{T/2}{f(t)g(t)}{}{t} \]
            この定義の下，次のベクトルの集合は $V$ の正規直行基底であることが知られている。
            （正規直交系であることは内積の計算で容易に確かめられる。 $V$ を張ることは数学の書籍に見出せる。）
            \[ \braces*{1/\sqrt{2},\cos\parens*{d\frac{t}{T}2\pi},\sin\parens*{d\frac{t}{T}2\pi}\;(d=1,2,3,\dots)} \]
            前述の Fourier 級数の定義では $\cos,\sin$ の項は基底を Fourier 係数で重み付けたものであるが，定数項だけは基底の $1/\sqrt{2}$ 倍に係数が掛かっており首尾一貫しない。
            ここを拘って係数と基底の積として表現すると $\inprod{f}{1/\sqrt{2}}/\sqrt{2} = \inprod{f}{1}/2 = a_0/2$ となるから $1/2$ が一般に使われているのだと思われる。
    \section{連続関数のFourier 級数の高周波成分は0に収束する}
        \label{連続関数のFourier 級数の高周波成分は0に収束する}
        \begin{shadebox}
            $f$は区間$[-\pi,\pi]$で連続とする。$f$のFourier係数を$c_k$とすると$\lim_{k\to\infty}c_k = 0$となる。
        \end{shadebox}
        \begin{proof}
            \quad\par
            まず
            \[ \lim_{k\to\infty}\integrate{-\pi}{\pi}{f(x)\sin kx}{}{x} = 0 \]
            を示す。ここで
            \[ \integrate{-\pi}{\pi}{f(x)\sin kx}{}{x} = \integrate{0}{\pi}{(f(x)-f(-x))\sin kx}{}{x} \]
            であり，$f(x)-f(-x)$は$[0,\pi]$で連続なので結局のところ区間$[0,\pi]$の任意の連続関数$f$に対して
            \[ \lim_{k\to\infty}\integrate{0}{\pi}{f(x)\sin kx}{}{x} = 0 \]
            が成り立つことを示せば良い。
            $f$は閉区間で連続だから有界であり，特に一様連続なので，任意の$\varepsilon>0$に対してある$\delta>0$が存在して
            $\forall (x,y\in[0,\pi],|x-y|\leq\delta),|f(x)-f(y)|<\varepsilon$が成り立つ。
            $k \geq 2,\pi/k < \delta$となるように$k$を十分大きくとる。
            \begin{align*}
                \integrate{0}{\pi}{f(x)\sin kx}{}{x} &= \sum_{l=0}^{\floor{k/2}-1}\left(\integrate{2l\pi/k}{(2l+1)\pi/k}{f(x) \sin kx}{}{x} + \integrate{(2l+1)\pi/k}{(2l+2)\pi/k}{f(x) \sin kx}{}{x}\right) \\
                &\quad + \integrate{2\floor{k/2}\pi/k}{\pi}{f(x)\sin kx}{}{x}
            \end{align*}
            最後の項は$k\to\infty$のとき$0$に収束する。
            最初の2項は$y=kx$と変数変換して
            \begin{align*}
                &\quad \sum_{l=0}^{\floor{k/2}-1}\frac{1}{k}\left(\integrate{2l\pi}{(2l+1)\pi}{f(y/k) \sin y}{}{y} + \integrate{(2l+1)\pi}{(2l+2)\pi}{f(y/k) \sin y}{}{y}\right) \\
                &= \frac{1}{k} \sum_{l=0}^{\floor{k/2}-1} \integrate{2l\pi}{(2l+1)\pi}{(f(y/k) - f(y/k+\pi/k)) \sin y}{}{y}
            \end{align*}
            となる。これの絶対値は
            \begin{align*}
                &\leq \frac{1}{k} \sum_{l=0}^{\floor{k/2}-1} \integrate{2l\pi}{(2l+1)\pi}{\abs{f(y/k) - f(y/k+\pi/k) \sin y}}{}{y} \\
                &< \frac{1}{k} \sum_{l=0}^{\floor{k/2}-1} \integrate{2l\pi}{(2l+1)\pi}{\varepsilon\abs{\sin y}}{}{y} = \frac{2\varepsilon}{k} \sum_{l=0}^{\floor{k/2}-1}1 = \frac{2\varepsilon}{k}\floor{k/2} \leq \frac{2\varepsilon}{k}\frac{k}{2} = \varepsilon
            \end{align*}
            同様にして
            \[ \lim_{k\to\infty}\integrate{-\pi}{\pi}{f(x)\cos kx}{}{x} = 0 \]
            も示せる。
        \end{proof}
    \section{\texorpdfstring{$\integrate{-\pi}{\pi}{\frac{\sin nx}{2\tan\frac{x}{2}}}{}{x} = \pi\;(n\in\naturalNumbers)$}{sin(nx)/(2 tan(x/2))の[-π,π]での積分はπ}}
        \label{integ{-pi}{pi}{(sin nx)/2tan(x/2)}{x} = pi}
        \begin{proof}
            \quad\par
            求めたい積分の値を$I_n$とする。
            分子の$\sin nx$について$nx = (n-1/2)x + x/2$として加法定理により展開すると
            \[ I_n = \frac{1}{2}\integrate{-\pi}{\pi}{\frac{\cos^2(x/2)}{\sin(x/2)}\sin(n-1/2)x}{}{x} + \frac{1}{2}\integrate{-\pi}{\pi}{\cos(x/2)\cos(n-1/2)x}{}{x} \]
            第一項の分子で$\cos^2(x/2) = 1-\sin^2(x/2)$とすると
            \begin{align*}
                I_n &= \frac{1}{2}\integrate{-\pi}{\pi}{\frac{\sin(n-1/2)x}{\sin(x/2)}}{}{x} - \frac{1}{2}\integrate{-\pi}{\pi}{\sin(x/2)\sin(n-1/2)x}{}{x} + \frac{1}{2}\integrate{-\pi}{\pi}{\cos(x/2)\cos(n-1/2)x}{}{x} \\
                &= \frac{1}{2}\integrate{-\pi}{\pi}{\frac{\sin(n-1/2)x}{\sin(x/2)}}{}{x} + \frac{1}{2}\integrate{-\pi}{\pi}{\cos nx}{}{x} \quad (\text{↑の第2,3項を加法定理でまとめた}) \\
                &= \frac{1}{2}\integrate{-\pi}{\pi}{\frac{\sin(n-1/2)x}{\sin(x/2)}}{}{x} \quad (\because\;\text{第2項は}0)
            \end{align*}
            さらに$(n-1/2)x = (n-1)x + x/2$として加法定理を使うと
            \[
                I_n = \frac{1}{2}\integrate{-\pi}{\pi}{\frac{\sin(n-1)x\;\cos(x/2)}{\sin(x/2)}}{}{x} + \frac{1}{2}\integrate{-\pi}{\pi}{\cos(n-1)x}{}{x} = I_{n-1} +
                \begin{cases}
                    \pi & (n=1) \\
                    0 & (\text{otherwise})
                \end{cases}
            \]
            これと$I_0 = 0$より定理が成り立つ。
        \end{proof}
    \section{Fourier 級数の応用}
        \subsection{\texorpdfstring{$\integrate{0}{\infty}{\frac{\sin x}{x}}{}{x} = \pi/2$}{sin(x)/xの0から∞までの積分はπ/2}}
        \begin{proof}
            \quad\par
            複素積分を使った導出がよく知られているが，実解析の範囲で導出してみる。
            この積分が収束することは，$(\sin x)/x$が$[0,\pi/2]$で連続であることから$[0,\pi/2]$で積分可能であること，および$\sin x = (-\cos x)'$を用いて部分積分することで$[\pi/2,\infty)$で積分が絶対収束することから示せる。
            よって，
            \begin{align*}
                \integrate{0}{\infty}{\frac{\sin x}{x}}{}{x} &= \lim_{n\to\infty}\integrate{0}{n\pi}{\frac{\sin x}{x}}{}{x} = \lim_{n\to\infty}\integrate{0}{\pi}{\frac{\sin nt}{t}}{}{t} \quad (t=x/n\text{とおいた}) \\
                &= \lim_{n\to\infty}\frac{1}{2}\integrate{-\pi}{\pi}{\frac{\sin nt}{t}}{}{t} = \lim_{n\to\infty}\frac{1}{2}\integrate{-\pi}{\pi}{\left(\frac{1}{t} - \frac{1}{2\tan(t/2)} + \frac{1}{2\tan(t/2)}\right)\sin nt}{}{t}
            \end{align*}
            $\frac{1}{t} - \frac{1}{2\tan(t/2)}$は$[-\pi,\pi]$で連続だから\ref{連続関数のFourier 級数の高周波成分は0に収束する}より
            \[ \lim_{n\to\infty}\frac{1}{2}\integrate{-\pi}{\pi}{\left(\frac{1}{t} - \frac{1}{2\tan(t/2)}\right)\sin nt}{}{t} = 0 \]
            であること，および\ref{integ{-pi}{pi}{(sin nx)/2tan(x/2)}{x} = pi}より証明が済む。
        \end{proof}